\section{Model Architecture}
\label{sec:model}

We instantiate the three-stage pipeline of
Section~\ref{sec:non_static} as illustrated in
Figure~\ref{fig:overview}. Stage~1 is a \emph{sequence encoder}
that maps a full activation tensor
$A \in \mathbb{R}^{L \times T \times D}$ to a sequence-level
uncertainty logit. Stage~2 freezes that encoder and applies it to a
sequence of centered token windows $\{A^{(t)}\}_{t=1}^{T}$, producing
per-window embeddings that a \emph{Mamba-2 temporal head}
\citep{albertgu2024} consumes to emit per-window logits
$\{\hat{p}^{(t)}\}_{t=1}^{T}$. A top-$k$ multiple-instance learning
(MIL) rule aggregates those logits into a sequence-level score, against
which the temporal head is supervised on per-atom segments. Stage~3
then refines the temporal head on whole passages: with the encoder
still frozen, the head is fine-tuned against per-atom auxiliary labels
obtained from atomic-fact verification, using a smoothed top-$k$
aggregation within each atom span. This shifts the head's loss surface
from scoring isolated atoms to ranking uncertainty \emph{across} the
atoms that co-occur in a passage. We describe the encoder first, the
temporal head second, and the passage-level refinement last.

\subsection{Stage~1: Encoder}
\label{sec:act_vit_subsec}

\paragraph{Overview.} The encoder consists of three components (see
Figure~\ref{fig:overview}, top): (1) a Pooling (Pool) layer that reduces
the spatial dimensions of the activation tensor (layers and tokens) to a
fixed, predefined size $(L_p, T_p)$; (2) a per-LLM Linear Adapter (LA)
that compresses and aligns the LLM-specific hidden dimensions into a
shared feature space; (3) a ViT-based backbone (ViT-B), which processes
the resulting tensor using a Vision Transformer architecture, enhanced
with shared positional encodings within each patch. A linear head on top
of the backbone emits a single sequence-level uncertainty logit. We
expand on each component below.

\paragraph{Pool.} Similarly to image resizing, we propose to reduce the
size of activation tensors before processing them with our architecture:
as a preprocessing step, we apply a max-pooling operation across the
spatial activation dimensions (see inset illustration). Specifically, we
consider two pooling approaches. The first maintains full input
resolution at the level of tokens and applies 1d pooling on the layer
dimension only; this approach is preferred when expecting relatively
concise LLM generations, or when the maximum generation length remains
computationally manageable. Alternatively, a second approach adopts a
more general 2d pooling on both spatial activation dimensions, allowing
us to more efficiently deal with potentially longer LLM responses. As we
will detail next, we mainly adopt the first approach and additionally
experiment with the second in Section~\ref{sec:experiments}.

In our applications, the number of generated tokens $T$ naturally varies
depending on the input prompt, and the number of LLM layers $L$ may vary
across models in $\mathcal{M}$. With guiding principle (i) in mind, it
becomes important to ensure our architecture can robustly deal with
both these varying parameters. Accordingly, we design the two pooling
operations so that, regardless of the input size, their outputs
$A_p \in \mathbb{R}^{L_p \times T_p \times D}$ always match a predefined
shape $(L_p, T_p)$ in their spatial dimensions (see
Algorithm~\ref{alg:pool}, in Appendix~\ref{app:pool}). This property
ensures the downstream ViT backbone is always fed with tensors of a
consistent spatial shape. Crucially, the same Pool module is reused
unchanged in Stage~2 with a smaller token target $W_p < T_p$, so a
single learned encoder operates on both full sequences and short
windows.

\paragraph{LA.} Following pooling, we map the feature dimension to a
shared hidden size via a linear transformation. We instantiate a
dedicated linear adapter for each LLM $M \in \mathcal{M}$, all
projecting to a shared output dimensionality $D'$:
\begin{equation}
\mathcal{L} = \left\{ LA_M : A_p \mapsto A_p W_M \mid W_M \in \mathbb{R}^{D_M \times D'},\ M \in \mathcal{M} \right\},
\end{equation}
where $D_M$ is the hidden dimension of model $M$ and $A_p$ denotes its
pooled activation tensor; the result has shape
$L_p \times T_p \times D'$. Applying a single linear transformation
$LA_M$ jointly across all layers and tokens of the activation tensor
for a given LLM $M$---rather than using separate projections per layer
or token---is motivated by the inherent alignment of feature spaces
within a single transformer: residual connections promote consistency
across layers, while transformer weight sharing across tokens ensures
a uniform representation structure along the token axis.

Resorting to LLM-specific LAs is particularly relevant in relation to
principle (i). It ensures that the downstream architectural components
process objects of consistent shape that are meaningfully comparable
across the internal representation spaces of distinct LLMs in
$\mathcal{M}$. Motivated by recent literature on representation
alignment \citep{huh2024,morcos2018,moschella2022,kornblith2019}, we
hypothesize that, despite differences in architectural design, training
corpora, weight initializations and learning schemes, the hidden
representations of LLMs can be put in approximate correspondence by
some linear transformations, which we aim to learn via the adapters in
$\mathcal{L}$. This contributes to rendering our multi-LLM setup not
only computationally possible, but also meaningful from a learning
perspective and, as we show in the next sections, beneficial in terms
of generalization performance. For an illustrative motivating example
of how a linear mapping can tackle misalignment in the representation
space, we refer the interested reader to Appendix~\ref{app:permutation},
where we specifically consider those induced by neuron symmetries, as
recently studied in \citep{navon2023,samuelkainsworth2022}.

\paragraph{ViT-B.} In view of the proposed analogy between activation
tensors and images, the output of any LA module is processed
\`a la ViT \citep{dosovitskiy2020}. It is first rearranged into a
sequence of non-overlapping activation patches of size $(p_H, p_W)$. If
$L_p$ and $T_p$ are divisible by $p_H$ and $p_W$, this gives
$H = L_p/p_H$ and $W = T_p/p_W$ patches; the resulting tensor is denoted
$A_{\text{patches}} \in \mathbb{R}^{(H \cdot W) \times p_H \times p_W \times D'}$.
We augment each patch with a per-patch learnable Positional Encoding
(PE), providing intra-patch ordering over (pooled) activation pixels.
The Transformer's native PEs provide a global ordering over activation
patches, and combined with our per-patch PEs, allow reconstructing each
pooled pixel's original position. Similarly to ViT, the enriched
activation patches are flattened, passed through a shared linear layer,
and processed by a stack of Transformer blocks. Overall, using a
vision-like backbone aligns with the analogy between activation tensors
and images and, as we will experimentally show, reveals to be an
effective inductive bias.

\paragraph{Stage~1 head and training.} A linear head on top of the
ViT-B backbone produces a single sequence-level uncertainty logit
$z \in \mathbb{R}$, trained with binary cross-entropy against the
sequence-level label $y$:
\begin{equation}
    \mathcal{L}_{\text{stage 1}} = \mathrm{BCE}\bigl(\sigma(z),\ y\bigr).
\end{equation}
We select the best checkpoint by validation AUROC. The learned features
just before the linear head form the per-sequence embedding that
Stage~2 reuses on a per-window basis.

\subsection{Stage~2: Mamba-2 Temporal Head with Top-$k$ MIL}
\label{sec:temporal_head}

Stage~2 keeps the Stage~1 encoder frozen and learns to localize
uncertainty over generated-token windows under sequence-level
supervision (see Figure~\ref{fig:overview}, bottom). It comprises
(1) a frozen window encoder reusing Stage~1, (2) a Mamba-2 temporal
head over per-window embeddings, and (3) a top-$k$ MIL aggregation that
maps per-window logits to a sequence-level score for training.

\paragraph{Frozen window encoding.} For each generated position
$t \in \{1, \dots, T\}$, we extract the windowed activation tensor
$A^{(t)} \in \mathbb{R}^{L \times W \times D}$ defined in
Section~\ref{sec:notation} as $W$ consecutive token columns of $A$
centered at $t$ (with edge padding near response boundaries). Each
$A^{(t)}$ is processed by the same Pool module used in Stage~1 with a
window-sized token target $(L_p, W_p)$, then passed through the frozen
encoder to produce a per-window embedding
$e^{(t)} \in \mathbb{R}^{d_{\text{model}}}$, taken from the
penultimate features of Stage~1's backbone (i.e., before the
sequence-level linear head). Stacking the $T$ windows yields a temporal
sequence
$E = (e^{(1)}, \dots, e^{(T)}) \in \mathbb{R}^{T \times d_{\text{model}}}$
on which the temporal head operates. Reusing the Stage~1 Pool with a
smaller target ensures that full-sequence and per-window encodings
share the same operation; for example, $(L, T)=(33, 200) \to (L_p, T_p)=(16, 100)$
in Stage~1 is paired with $(L, W)=(33, 6) \to (L_p, W_p)=(16, 3)$ in
Stage~2.

\paragraph{Mamba-2 temporal head.} A residual stack of Mamba-2 blocks
\citep{albertgu2024} interleaved with LayerNorm processes $E$ along the
token axis, followed by a linear head that emits one logit per
position:
\begin{equation}
    \hat{z}^{(t)} = h\bigl(\mathrm{Mamba\text{-}2}(E)_t\bigr)
    \quad \text{for } t = 1, \dots, T.
\end{equation}
We adopt Mamba-2 over a Transformer here because activation sequences
in long-form generation can extend into the thousands of tokens; its
linear-time selective state-space recurrence keeps Stage~2 viable for
real-time, streaming use, where each new generated token requires only
a constant-memory update of the temporal head. The corresponding
per-window uncertainty score is
$\hat{p}^{(t)} = \sigma(\hat{z}^{(t)}) \in [0, 1]$ and constitutes our
\emph{localization output}.

\paragraph{Top-$k$ MIL aggregation.} Because no token- or window-level
labels are available, we train the temporal head from sequence-level
supervision via top-$k$ multiple-instance learning. After masking
padded positions, we select the $k$ largest per-window logits per
sequence and average them into a single sequence-level logit:
\begin{equation}
    \hat{z} = \frac{1}{k} \sum_{t \in \mathrm{Top}_k\bigl(\{\hat{z}^{(t)}\}_t\bigr)} \hat{z}^{(t)},
    \qquad
    \hat{p} = \sigma(\hat{z}),
\end{equation}
where $\mathrm{Top}_k(\cdot)$ denotes the indices of the $k$ largest
unmasked logits. Stage~2 is then trained with binary cross-entropy
against the same sequence label as Stage~1:
\begin{equation}
    \mathcal{L}_{\text{stage 2}} = \mathrm{BCE}(\hat{p},\ y).
\end{equation}
Mean-of-top-$k$ is more stable than alternative aggregations
(e.g., max or full mean) for long sequences: max is brittle to single
high-logit windows, while a full mean dilutes the few suspicious
regions with the majority of correct ones. Top-$k$ implements the
intuition that an uncertain or incorrect response typically contains
only a few problematic spans, and lets the model discover those spans
under weak supervision.

\subsection{Stage~3: Passage-level Refinement with Per-atom Supervision}
\label{sec:stage3}

Stages~1 and~2 train against \emph{sequence-level} labels on per-atom
segments: at training time the temporal head sees one atom at a time,
and its loss reflects how well it scores that atom in isolation. In
deployment, however, the head must operate over a whole passage whose
atoms co-occur and compete for attention; a head tuned only on isolated
atoms tends to assign similar scores to every region of a passage and
fails to localize the uncertain few. Stage~3 addresses this by
fine-tuning the temporal head on whole passages under \emph{per-atom}
auxiliary supervision, while keeping the encoder frozen.

\paragraph{Inputs.} The training unit is now a full passage activation
tensor $A \in \mathbb{R}^{L \times T \times D}$ paired with $M$ atom
spans $\{(\ell_m, r_m)\}_{m=1}^{M}$ obtained from an atomic-fact
verification pipeline, where the $m$-th atom occupies generated-token
positions $\ell_m, \ldots, r_m$. Each span carries an auxiliary
hallucination label $\tilde{y}_m \in \{0, 1\}$ broadcast from the
fact-level verification result. As discussed in
Section~\ref{sec:non_static}, these labels are noisy -- the span may
contain non-hallucinated tokens, and the relevant activation signal can
even fall outside the span -- so we treat them as coarse auxiliary
supervision rather than as ground-truth token labels.

\paragraph{From per-token logits to per-atom scores.} The frozen
encoder and the Stage~2 temporal head jointly produce a per-token logit
sequence $\hat{z}^{(1)}, \ldots, \hat{z}^{(T)}$ over the passage. We
convert these to per-atom scores in two steps. First, we smooth the
logit sequence along the token axis with a centered length-$W$ window:
\begin{equation}
\bar{z}^{(t)}
= \sum_{u \in \mathcal{W}(t)} \alpha_u^{(t)} \, \hat{z}^{(u)},
\qquad
\alpha_u^{(t)} \propto \exp(\hat{z}^{(u)}/\tau),
\label{eq:soft-topk-smooth}
\end{equation}
where $\mathcal{W}(t)$ is the centered length-$W$ window at position
$t$ (with masking near response boundaries) and $\tau > 0$ controls the
sharpness: $\tau \to 0$ recovers the per-window max, while $\tau \to
\infty$ recovers the arithmetic mean. The smoothing absorbs the
window-level uncertainty signal already produced by Stage~2 in a
differentiable way, so gradients flow back into the head from any
informative position within the window. Second, we aggregate the
smoothed logits within each atom span $[\ell_m, r_m]$. By default we
use a top-$k$ mean,
\begin{equation}
\hat{z}_m
= \frac{1}{k_m}
\sum_{t \in \mathrm{Top}_{k_m}\!\bigl(\{\bar{z}^{(t)}\}_{t = \ell_m}^{r_m}\bigr)}
\bar{z}^{(t)},
\qquad
k_m = \max\!\bigl(k_{\min}, \lceil \kappa \cdot (r_m - \ell_m + 1) \rceil\bigr),
\label{eq:topk-segment}
\end{equation}
with $\kappa \in (0, 1]$ controlling the fraction of in-span tokens
that participate in the atom's score. Setting $\kappa = 1$ averages
over every token in the span; smaller $\kappa$ peaks the gradient on
the most uncertain tokens within the span and is useful for long
atoms. A softmax-weighted in-span max
($\hat{z}_m \propto \sum_t \mathrm{softmax}(\bar{z}^{(t)}/\tau_{\text{seg}}) \,
\bar{z}^{(t)}$) is available as an ablation.

\paragraph{Stage~3 loss.} The temporal head is trained with a per-atom
binary cross-entropy against the auxiliary labels:
\begin{equation}
\mathcal{L}_{\text{stage 3}}
= \frac{1}{\sum_m \mathbb{1}[\text{atom } m \text{ valid}]}
\sum_m \mathbb{1}[\text{atom } m \text{ valid}] \cdot
\mathrm{BCE}\bigl(\sigma(\hat{z}_m),\ \tilde{y}_m\bigr).
\end{equation}
Compared with the sequence-level MIL objective of Stage~2, this loss
forces the head to reason \emph{within} the passage: every atom
contributes its own term, so the head cannot trivially explain a
passage label by saturating a single window -- it must rank atoms
against each other under the same encoder.

\paragraph{Inference.} At inference, a single forward pass through the
Stage~3 detector yields the per-token logits over the whole passage,
from which we read off both outputs of our problem statement
(Section~\ref{sec:notation}): the per-window scores
$\{\hat{p}^{(t)}\}_{t=1}^{T} = \{\sigma(\hat{z}^{(t)})\}_{t=1}^{T}$ used
to localize the generated-token regions whose hidden-state activations
most contribute to the prediction, and a sequence-level score obtained
by reusing the Stage~2 top-$k$ MIL aggregation over the same per-token
logits. The per-window scores can be used directly for visualization or
as inputs to downstream verifiers and correction routines on long-form
responses; when atom spans are available, the Stage~3 per-atom scores
$\{\sigma(\hat{z}_m)\}_{m=1}^{M}$ provide a higher-level summary
aligned with claim-level verification.
